\section{Illustration: the Philippine coupled run}
\label{sec:results}

We illustrate the implemented forward pass on the Philippine calibration ($\Mind=8$
industries, $\Igood=5$ goods; electricity is industry $m_e=2$), coupled to a
\CLEWS{}/\OSeMOSYS{} scenario pair. The application tests whether the principal forward
mappings are accepted by the economy solve and whether their responses are consistent
with the mechanisms of \Cref{sec:channels}; it is not a policy evaluation. The coupled
application runs on the Philippine v16 pair --- a policy-composition reform that
includes a coal moratorium from 2027, a binding limit on offshore wind, and
land-conversion carbon accounting --- and its quoted magnitudes derive from the
archived run record (\Cref{app:run-record}) rather than being transcribed by hand. The
composition analysis of \Cref{sec:decomposition} was established on the earlier v12
pair; every table and figure names its evidence base. Rates, wages, and quantities are
reported as percent changes from baseline, and all magnitudes are relative changes
(\Cref{sec:calibration}).

\subsection{The coupled application}
\label{sec:coupled}
The \chan{coupled} experiment applies, in one reform: the \CLEWS{} electricity-cost
signal --- a levelized cost, \Cref{sec:signal} --- through the composite transmission
of \Cref{sec:routes} (inter-industry cost-push plus a recycled household wedge); the
public component of grid investment; and the Global Burden of Disease health mapping.
Alongside the equilibrium return and activity paths it emits a \$50/tCO\textsubscript{2}
carbon-price input for the energy model. The emitted inputs are inert in a single pass:
the underlying \CLEWS{} reform was solved without an emissions penalty, so no
carbon-price effect enters any number reported here.

The realized inputs, from the run record: the mean household electricity-price wedge is
$+\cplWedgePct\%$; the public grid-investment delta is $\cplInvestCumPctGDP\%$ of GDP
cumulatively --- effectively zero; and fine-particulate exposure changes by
$\cplPMPct\%$, which the health mapping converts to roughly \cplDeaths{} avoided deaths
(on a GBD age profile) and a working-time gain of about \cplMorbFTE{}
full-time-equivalent workers per year.

The response is a contraction in both the transition window and the long run
(\Cref{tab:phl-coupled}): steady-state output is $\cplYss\%$ relative to baseline, and
the \cplWindow{} average is $\cplYwin\%$. The adjustment is factor-price led: the wage
carries most of it ($\cplWss\%$ at the steady state), aggregate labour is close to
unchanged ($\cplLss\%$), and the capital stock declines with output ($\cplKss\%$).
\Cref{fig:transition} traces the path: the contraction deepens through the 2030s and
2040s, reaches its largest output deviation (about $-0.7\%$) near 2050, and settles
toward the steady state.

One qualification belongs beside the headline rather than in a footnote. Nearly all of
this response transmits through the inter-industry cost-push leg, which maps the
sectoral electricity-cost change into unit costs as if it were an exogenous input-price
shift; it is a reduced-form mapping, not a factor-demand system. A structural
alternative that routes the same cost signal through sectoral productivity reverses the
sign of the aggregate response ($\atrTfpY\%$ against $\atrCostPushY\%$ for the
cost-push leg alone, on the v12 pair; \Cref{sec:decomposition}). The headline is
therefore a joint statement about the scenario and the transmission assumption, and it
should not travel without the transmission named.

\begin{table}[t]
\centering\footnotesize
\caption{The Philippine coupled run (percent change from baseline; transition-window
average and \OGCore{} steady state). Evidence base: \cplBase. Source: the archived run
record (\Cref{app:run-record}).}
\label{tab:phl-coupled}
% GENERATED — see numbers.tex header.
% Evidence base: the Philippine application (PEP vs Base), OG-PHL M=8
\begin{tabular}{@{}lrr@{}}
\toprule
 & 2025--2034 average & steady state \\
\midrule
output $Y$        & $-0.279$ & $-1.001$ \\
consumption $C$   & $-0.127$ & $-0.720$ \\
capital $K$       & $-0.303$ & $-0.950$ \\
labour $L$        & $0.007$ & $0.025$ \\
return $r$        & $-0.005$ & $-0.192$ \\
wage $w$          & $-0.244$ & $-0.929$ \\
\bottomrule
\end{tabular}

\end{table}

\begin{figure}[t]
\centering
\includegraphics[width=0.82\linewidth]{v16_macro_transition}
\caption{The coupled transition path (percent change from baseline): output,
consumption, capital, and labour, 2025--2074. The largest output deviation is about
$-0.7\%$ near 2050. Evidence base: \cplBase; from the run record's figure set.}
\label{fig:transition}
\end{figure}

\subsection{Composition of the coupled response}
\label{sec:decomposition}
Within the applied configuration the composition is largely fixed by construction: the
carbon input is emitted only, the public-investment delta is near zero, and the health
mapping's aggregate contribution is below reporting precision. The coupled response
must therefore coincide with the response to the energy composite alone, up to those
near-zero terms --- and a matched pair of runs on the v12 scenario pair confirms that
it does (\Cref{tab:attribution}): the coupled steady state ($\atrCoupledY\%$) equals
the energy-composite steady state ($\atrCompositeY\%$) at printed precision on output,
consumption, capital, and labour, while the wage differs in the fourth decimal --- the
two are not bit-identical. This is an accounting identity of the configuration, not a
finding that channels do not interact.

The instructive split is inside the composite. Run separately on the same pair, the
inter-industry cost-push leg accounts for $\atrCostPushY\%$ of the $\atrCoupledY\%$
total and the household wedge for $\atrWedgeY\%$: firms' unit costs, not household
electricity bills, carry the aggregate effect. The table's final row repeats the
transmission sensitivity stated above: the structural productivity routing of the same
signal yields $\atrTfpY\%$.

Two reading rules follow. First, the standalone carbon experiment reported in the run
battery applies a consumption-side carbon tax inside the economy model --- a treatment
the coupled configuration deliberately omits --- so its response must not be summed
with, or netted against, the coupled headline. Second, these composition results are
established on the v12 pair; the matched runs on the illustration's own base are
queued, and until they exist the split travels with its base named.

\begin{table}[t]
\centering\footnotesize
\caption{Composition of the coupled steady-state response (percent change from
baseline). Evidence base: \atrBase; re-verification on the illustration's base is
pending. Source: the archived run record (\Cref{app:run-record}).}
\label{tab:attribution}
% GENERATED — see numbers.tex header. Composition battery pending.
\begin{tabular}{@{}c@{}}
\genPENDING{matched composition battery on the current calibration}\\
\end{tabular}

\end{table}

\subsection{Distributional incidence}
\label{sec:incidence}
The age- and ability-structured household side reports who bears the adjustment. On
the v16 run the consumption-equivalent incidence across lifetime-income groups is
strongly heterogeneous and not monotone in income (\Cref{fig:incidence}): the middle
groups gain (about $+2.9\%$ for the 50th--70th percentile group), the upper-middle
groups bear the largest losses (about $-1.6\%$ for the 90th--99th), and the poorest
and richest groups are close to unchanged. This is not a budget-share pattern:
baseline energy spending is roughly $1.4\%$ of the consumption basket in every group
in this calibration, so the household wedge itself is nearly proportional. The
dispersion instead runs through factor incomes and the fiscal closure --- which
ability types earn the falling wage, hold the shrinking capital stock, and receive the
recycled revenue and adjusting transfers. A representative-household core would report
only the aggregate ($\cplCss\%$); the incidence statement is what the age--ability
structure adds, and it is the layer a distribution-blind coupling cannot produce.

\begin{figure}[t]
\centering
\includegraphics[width=\linewidth]{v16_incidence}
\caption{Distributional incidence of the coupled run: consumption-equivalent change by
lifetime-income group; the middle panel shows baseline energy budget shares are nearly
uniform ($\approx 1.4\%$), so the pattern is not budget-share driven. Evidence base:
\cplBase; from the run record's figure set.}
\label{fig:incidence}
\end{figure}

\subsection{The capital pair}
\label{sec:capital-results}
The \chan{capital\_intensity} channel is decomposed exactly via the firm identity of
\Cref{sec:ch-capint} (\Cref{tab:capital-pair}, measured on an earlier steady-state solve
of the lever): raising the energy capital exponent
$12.2\%$ collapses the electricity price $42.8\%$; demand is elastic enough that output
\emph{rises} $20.9\%$, yet capital \emph{falls} $22.4\%$ because the price collapse shrinks
revenue---and the identity $\hat K_E=\hat\gamma_E\hat p_E\hat Y_E$ reproduces that response
to within $0.03\%$. It is a factor-share and energy-price lever, not a capital-draw-in.

\subsection{Validation}
\label{sec:validation}
Validation is layered because the tests establish different claims
(\Cref{tab:validation-layers}). The grouped run battery executes the link in its own
environment and the country \OGCore{} model in the country model's environment,
exchanging data files rather than sharing an interpreter. All battery items pass the
solve-completion gate. Selected macro aggregates of the \emph{v9} scenario pair are
regression-locked in \texttt{results/golden.json} and re-checked automatically; the
v12 battery and the v16 illustration are not yet regression-locked.

\begin{table}[t]
\centering\footnotesize
\caption{Validation layers and their current boundaries.}
\label{tab:validation-layers}
\begin{tabularx}{\linewidth}{@{}>{\bfseries\raggedright\arraybackslash}p{0.20\linewidth} L L@{}}
\toprule
Evidence & What it establishes & What it does not establish \\
\midrule
Transform and guard tests & dimensions, bounds, concordances, and failure behaviour & economic calibration or equilibrium response \\
Controlled mechanism runs & solve completion, response direction, and selected identities & joint attribution or application magnitude \\
Golden aggregate record & reproducibility of selected macroeconomic outputs (v9 pair) & every sectoral or emitted interface artifact; later scenario pairs \\
Coupled forward run & joint solution of the price, public-investment, and health mappings & closed-loop equilibrium \\
Matched composite pair & the coupled response equals its energy composite up to the near-zero channels (v12 pair) & the same identity on other scenario pairs or bases \\
\CLEWS{} re-solve test & a controlled demand patch survives copy, solve, and read-back & passage of emitted return and activity paths end to end \\
Closed-loop iteration & not yet tested & convergence and propagated-error control \\
\bottomrule
\end{tabularx}
\end{table}

One documented exception is that the lives-saved health reform leaves an intrinsic
Walras residual that scales
with the target, so its steady-state resource-constraint gate is loosened to $10^{-5}$
(an order of magnitude under \OGCore{}'s own transition-path default), while the
deaths-added direction stays at the tight default.

\caveat{The electricity-cost input is a scenario-derived levelized cost, not a marginal
price: the commodity-balance dual is degenerate in this case, so no marginal reading is
available or implied. All magnitudes are relative changes --- the unit/deflator bridge
is not yet audited (\Cref{sec:calibration}), so currency-level readings of carbon and
investment are illustrative. The morbidity response is a placeholder calibration
pending dose--response data; the mortality target, age profile, and morbidity profile
are GBD-sourced.}
